Model calibration served as a critical validation step to ensure that the simulation accurately reproduced the current real-world state of the Municipality of Bacolod prior to the introduction of optimization policies. This process necessitated reconciling the discrepancies observed between micro-level data gathered through barangay interviews and the macro-level reality established by the Municipal Environment and Natural Resources Office audit.

\subsection{Addressing Reporting Bias in Input Data}
A significant challenge identified during the initialization phase was the marked divergence between self-reported compliance rates and the observed municipal average.  Aggregating the compliance rates claimed by barangay officials yielded a calculated municipal average ranging from approximately 40\% to 50\%, with specific localities such as Binuni and Demologan reporting rates as high as 95\% and 60\%, respectively. However, the municipal-wide audit conducted by the MENRO indicated an actual effective segregation rate of approximately 10\%. This discrepancy quantitatively captured the critical ``Act vs. Reality'' gap described by Yazawa et al. \cite{Yazawa2025}, where reported compliance at the barangay level frequently deviated significantly from the actual ground-level implementation of RA 9003.

This discrepancy suggested the presence of social desirability bias in the interview data, where local officials may have overstated compliance to present a favorable political image. It also corroborated the findings of Abordo and Dalugdog \cite{Abordo2025}, who observed that high levels of reported awareness or ``policy knowledge'' often failed to translate into consistent, proper segregation practice. Consequently, the simulation utilized the MENRO audit of 10\% as the ground truth target for calibration.

\subsection{Calibration of Initial Compliance States}
A critical divergence was observed between the compliance rates self-reported by barangay officials during Key Informant Interviews (Appendix B) and the objective audit data provided by the Municipal Environment and Natural Resources Office (MENRO). While aggregate interview data suggested a municipal compliance rate exceeding 50\%, the MENRO’s waste characterization study confirmed a functional segregation rate of only 10\% to 15\%.

To ensure the Agent-Based Model reflected the true ``Status Quo,'' the initialization parameters in \texttt{barangay\_config.py} were calibrated to match the MENRO ground truth rather than the optimistic self-reports. This decision rested on three specific analytical pivots:

\begin{table}[htbp]
    \centering
    \caption{Comparison of Acquired Compliance Data vs. Calibrated Simulation Initialization}
    \label{tab:calibration_data}
    \small
    \begin{tabular}{l c c l}
        \toprule
        {Barangay} & {Acquired Data} & {Calibrated $S_0$} & {Primary Calibration Factor} \\
        & \textit{(Self-Reported)} & \textit{(Simulation Start)} & \\
        \midrule
        Brgy Babalaya     & 100\% & 14\% & Adjusted for lack of funding/monitoring \\
        Brgy Binuni       & 95\%  & 15\% & Decoupled awareness from actual practice \\
        Brgy Mati         & 70\%  & 11\% & Adjusted for high habit decay rate \\
        Brgy Liangan East & 65\%  & 14\% & Removed temporary project-based spikes \\
        Brgy Demologan    & 60\%  & 11\% & Standardized to municipal baseline \\
        Brgy Ezperanza    & 20\%  & 14\% & Aligned with middle-class control group \\
        Brgy Poblacion    & 2\%   & 2\%  & {Anchor Point} (Aligned with Audit) \\
        \bottomrule
    \end{tabular}
    \vspace{1ex}
    
    \raggedright
    \footnotesize
    \textit{Note: "Acquired Data" refers to unverified compliance rates reported during Key Informant Interviews (Appendix B). "Calibrated $S_0$" refers to the initialized compliance state in the Agent-Based Model, tuned to match the aggregate municipal segregation rate of $\approx$12\% verified by MENRO.}
\end{table}

Barangay Binuni officials reported a near-perfect compliance rate of 95\% (Appendix B.3). However, initializing the simulation at this level resulted in immediate global compliance exceeding 40\%, rendering the RL agent’s optimization tasks trivial. The simulation logic posited that the reported 95\% reflected \textit{awareness} rather than \textit{practice}. This distinction was critical; as noted by \cite{Paigalan2025}, positive attitudes (Mean=3.17) often coexisted with poor practices like open burning (Mean=2.90), creating an ``Intention-Action Gap'' that the model had to explicitly correct for.

Barangay Liangan East reported 60\% compliance driven by an Eco-brick initiative. However, longitudinal analysis suggested that project-based compliance was transient. To model the ``Status Quo'' before the DRL agent's intervention, the system initialized Liangan East at 14\%. This reflected the baseline habituated behavior of the population when specific project funds were not actively being disbursed, preventing the simulation from starting at an artificial peak.

The only data point where self-reporting matched the MENRO audit was Barangay Poblacion, which reported a 2\% compliance rate. This convergence served as the \textit{Anchor Point} for the simulation. Because the qualitative data and quantitative audit agreed on the lower bound, the simulation initialized Poblacion at 2\%. This validated the ``Urban Resource Trap'' theory, where high population density and anonymity dissolved social pressure, creating a distinct legislative challenge that the DRL agent had to solve through enforcement rather than social norms. This low baseline confirmed the assessment by Villanueva et al. \cite{Villanueva2021} that without strict monitoring mechanisms, the ``status quo'' in high-density areas defaulted to non-compliance despite the existence of legal frameworks.

By calibrating the initial states to this conservative baseline (2\% to 15\%), the simulation stabilized at a Global Average Compliance of $\approx$12.5\% at $t=0$. This provided a valid, non-inflated baseline against which the DRL agent's policy improvements could be accurately measured.
